\documentclass[10pt,a4paper]{amsart}
\usepackage[margin=19mm]{geometry}
\usepackage{amsmath,amssymb,mathtools}
\usepackage[scr=rsfs,cal=euler]{mathalfa}
\usepackage{xfrac}
\usepackage[unicode,hidelinks,bookmarks=false]{hyperref}
\newcommand{\ie}{i.e.\ }
\newcommand{\cf}{cf.\ }
\newcommand{\resp}{respectively\ }
\newcommand{\Subsection}{\subsection}
\providecommand{\nobreakdash}{\nobreak\hbox{-}}
\setlength{\emergencystretch}{2em}
\multlinegap0pt
\usepackage{amssymb}
\usepackage{float}
\newtheorem{lemma}{Lemma}
\def\M{\mathfrak{M}}
\title{Solving the Sylvester equation in Banach modules}
\author{Bogdan Djordjevi\'c}
\date{Source: arXiv:2605.13419v1 (CC BY 4.0)}
\begin{document}
\maketitle

\noindent\emph{Source and licence.} Solving the Sylvester equation in Banach modules. Version arXiv:2605.13419v1 (CC BY 4.0), distributed by arXiv under the Creative Commons Attribution 4.0 International licence, http://creativecommons.org/licenses/by/4.0/, as recorded in the arXiv licence metadata for this exact version and shown on the abstract page https://arxiv.org/abs/2605.13419v1. Full licence notice: \texttt{licenses/arxiv-2605.13419-CC-BY-4.0.txt}.\par\smallskip

\noindent\textit{Editorial scope.} Counted unit: the article's lemma uvsyl (source0.tex lines 632--645). The article's complete original proof of it is delivered with it (source0.tex lines 647--696, one \texttt{proof} environment of 50 lines). No mathematics is written, repaired or re-derived: the statement and the proof are the article's own lines, in the article's own notation, and no retained line is substituted or re-flowed.  Retained uncounted as the source-local context the proof needs: lines 564--566; lines 568--570; lines 571--573; lines 574--576; lines 577--579; lines 581--590.  Nothing was omitted from any retained span, so the package carries no omission record.  No retained line contains a citation, so the wrapper carries no bibliography and no bibliography entry.  No retained line contains a figure environment, an image-inclusion command or drawing code, so no image is delivered and the package declares no asset dependency.  Editorial formatting only, in addition: an AMS \texttt{amsart} preamble that restates the article's own declarations and macros used by the retained lines, namely the environments \texttt{lemma} on separate counters, together with the standard packages those lines need.  The retained environments print in this document as Lemma 1 in order of appearance, whereas the article prints the same items with the numbering of its own full document.  The complete native source of the article as distributed in the arXiv e-print for this exact version is bundled unchanged as \texttt{sources/200503/source0.tex}.
\begin{equation}\label{r}
r:=a^{-1}\bar{c}b+a\bar{c}b^{-1}.
\end{equation}

\begin{equation}\label{gensyl}
av+ub=\bar{c}.
\end{equation}

\begin{equation}\label{gensylvu}
au+vb=\bar{c}+a^{-1}\bar{c}b+a\bar{c}b^{-1}.
\end{equation}

 \begin{equation}\label{u+v}
u+v=a^{-1}cb^{-1}.
\end{equation}

 \begin{equation}\label{RajUchiuv}
a^3v+a^2vb+ub^3+aub^2=0.
\end{equation}

\begin{lemma}\label{2of4} Let $a$, $b$, $c$, and $\bar{c}$ be provided as above. Let $u$ and $v\in\M$ be arbitrary. Then the chain of equivalences is true for $u$ and $v$:
\begin{eqnarray}
\label{equivalence}
\begin{aligned}
&\left(\eqref{gensyl}\wedge\eqref{gensylvu}\right)\Leftrightarrow\left(\eqref{u+v}\wedge\eqref{gensylvu}\right)\Leftrightarrow\left(\eqref{gensyl}\wedge\eqref{u+v}\right)\\
&\Leftrightarrow\left(\eqref{RajUchiuv}\wedge\eqref{u+v}\right)\Leftrightarrow\left(\eqref{gensyl}\wedge\eqref{RajUchiuv}\right).
\end{aligned}
\end{eqnarray} 
In other words, if $u$ and $v$ solve any two equations enlisted in \eqref{equivalence}, then they solve all of them.
\end{lemma}

 \begin{lemma}\label{uvsyl}
Let $a$, $b$, $c$, and $\bar{c}$ be appropriately given elements, and let $r$ be provided as in \eqref{r}. For arbtirary two elements $u,v\in\M$ the following statements are equivalent: 
\begin{itemize}
\item[(a)] The ordered pair $(u,v)$ satisfies any two conditions enlisted in \eqref{equivalence}.
\item[(b)] The element $u$ satisfies
\begin{equation}
\label{auub}
au-ub=a\bar{c}b^{-1}\end{equation}
\item[(c)] The element $v$ satisfies:
\begin{equation}\label{avvb}
av-vb=-a^{-1}\bar{c}b.
\end{equation}
\end{itemize}
\end{lemma}

\begin{proof} We show the equivalence $(a)\Longleftrightarrow(b)$. The equivalence $(a)\Longleftrightarrow(c)$ is analogous.

$(b)\Rightarrow(a):$ Assume that \eqref{auub} holds. Then
$$\begin{aligned}
&a^2u-ub^2=a(au-ub)+(au-ub)b=\\
=&a\left(\int_0^{+\infty} e^{-at}acb^{-1}e^{-bt}dt\right)+\left(\int_0^{+\infty} e^{-at}acb^{-1}e^{-bt}dt\right)b=\\
=&\int_0^{+\infty}e^{-at}(a^2cb^{-1}+ac)e^{-bt}dt,
\end{aligned}$$
therefore
\begin{eqnarray}
\label{b2}
\begin{aligned}
&aub^{-1}-a^{-1}ub=a^{-1}\left(a^2u-ub^2\right)b^{-1}=\int_0^{+\infty}e^{-at}(acb^{-2}+cb ^{-1})e^{-bt}dt=\\
=&\int_0^{+\infty}e^{-at}\left(a\left(cb^{-2}\right)+\left(cb ^{-2}\right)b\right)e^{-bt}dt=cb^{-2}.\end{aligned}\end{eqnarray}
Since $r$ is the unique solution to $ar+rb=a^{-1}cb+acb^{-1}$, it is easy to see that
$$cb^{-1}+a^{-1}\bar{c}b-\bar{c}=r.$$
However, this implies that 
$$cb^{-2}=rb^{-1}+\bar{c}b^{-1}+a^{-1}\bar{c}.$$
Substituting the latter into \eqref{b2} gives
\begin{equation}
\label{invdiff}
aub^{-1}-a^{-1}ub=\bar{c}b^{-1}+rb^{-1}-a^{-1}\bar{c},\end{equation}
or, equivalently,
$$a^{-1}\bar{c}-a^{-1}ub=\bar{c}b^{-1}+rb^{-1}-aub^{-1}.$$
Set 
$$v:=a^{-1}\bar{c}-a^{-1}ub\equiv \bar{c}b^{-1}+rb^{-1}-aub^{-1}.$$ 
Then $av+ub=\bar{c}$ and $au+vb=\bar{c}+r$. By Lemma \ref{2of4} the pair $(u,v)$ satisfies all the conditions in \eqref{equivalence}.\\


$(a)\Rightarrow(b):$ Let $u$ and $v$ in $\M$ be provided such that \eqref{gensyl} and \eqref{RajUchiuv} hold. Denote by   $d:=au - ub$. Then $au = ub + d$, $aub^2 = ub^3 + db^2$, and $ub^3 = aub^2 - db^2$.
On the other hand, we have $av = \bar{c} - ub$, $
a^2 vb = a\bar{c}b - aub^2$, and $a^3 v = a^2 \bar{c} - a^2 ub$. Therefore
$$\begin{aligned}
0=&a^3 v + a^2 vb + aub^2 + ub^3
=\\
=&a^3 v + a^2 vb + 2aub^2 - db^2
=\\
=&a^2 \bar{c} - a^2 ub + a\bar{c}b - aub^2 + 2aub^2 - db^2=\\
=&a(a\bar{c} + \bar{c}b) + aub^2 - a^2 ub - db^2=\\
=&ac - \left(a(au - ub)b + db^2\right)=ac -\left(adb + db^2\right)=\\
=&ac - (ad + db)b.
\end{aligned}
$$
This is equivalent to
\[
acb^{-1} = ad + db
\Longleftrightarrow
au - ub = \int_0^\infty e^{-at}\, a c b^{-1}\, e^{-bt}\, dt.
\] 
\end{proof}

\end{document}
